\chapter{Automa\c{c}\~ao DNS e Web}
\label{cap3}

\section{Enquadramento do cap\'{\i}tulo}

Este cap\'{\i}tulo apresenta a implementa\c{c}\~ao dos pontos do enunciado associados \`a administra\c{c}\~ao de DNS e de servi\c{c}os Web, com foco na automatiza\c{c}\~ao de tarefas operacionais que, quando feitas manualmente, s\~ao mais demoradas e propensas a erro.

Do ponto de vista funcional, o DNS (\textit{Domain Name System}) garante a resolu\c{c}\~ao de nomes para endere\c{c}os IP. Isso permite que os servi\c{c}os sejam acedidos por dom\'{\i}nios leg\'{\i}veis em vez de endere\c{c}os num\'ericos. J\'a o servi\c{c}o Web, neste caso com Apache (\texttt{httpd}), responde aos pedidos HTTP e entrega o conte\'udo definido para cada dom\'{\i}nio atrav\'es de \textit{VirtualHosts}.

No contexto do projeto, estas duas camadas est\~ao diretamente ligadas: primeiro o DNS aponta o dom\'{\i}nio para o servidor, depois o Apache seleciona o \textit{VirtualHost} correto para servir a p\'agina esperada. Por isso, a integra\c{c}\~ao DNS+Web \'{e} tratada aqui de forma conjunta, com valida\c{c}\~ao em ambiente cliente-servidor.

\section{Quest\~ao 1: cria\c{c}\~ao autom\'atica de zona master DNS}
\label{sec:q1}

\subsection{Objetivo}

Na Quest\~ao 1, pretende-se automatizar a cria\c{c}\~ao de uma zona \textit{master} no servidor DNS, a partir da introdu\c{c}\~ao, por parte do utilizador, do nome do dom\'{\i}nio e do endere\c{c}o IP do servidor. O objetivo principal consiste em eliminar a necessidade de configura\c{c}\~ao manual no BIND, assegurando que a nova zona \'{e} criada de forma autom\'atica e que inclui, pelo menos, os registos do tipo \texttt{A} necess\'arios para a resolu\c{c}\~ao do dom\'{\i}nio para o IP do servidor DNS/Web.

Esta automatiza\c{c}\~ao permite reduzir erros de configura\c{c}\~ao, tornar o processo mais r\'apido e garantir uma estrutura uniforme na cria\c{c}\~ao de novas zonas DNS.

\subsection{Prepara\c{c}\~ao do servi\c{c}o DNS}

Antes da execu\c{c}\~ao do script, foi necess\'ario preparar a m\'aquina servidor com os pacotes e servi\c{c}os base do BIND, de forma a disponibilizar o ambiente necess\'ario \`a cria\c{c}\~ao e valida\c{c}\~ao das zonas DNS.

\begin{lstlisting}[language=bash,caption={Preparacao inicial do servico DNS}]
yum install bind bind-utils -y
systemctl enable named
systemctl start named
systemctl status named
\end{lstlisting}

Foram instalados os pacotes \texttt{bind} e \texttt{bind-utils}, essenciais para o servi\c{c}o DNS e para a valida\c{c}\~ao de configura\c{c}\~oes. Em seguida, o servi\c{c}o \texttt{named} foi ativado e iniciado, garantindo que o servidor estava operacional antes da execu\c{c}\~ao do script.

Ap\'os a instala\c{c}\~ao e ativa\c{c}\~ao do servi\c{c}o, foi criada a diretoria de trabalho do projeto, onde foram armazenados os scripts desenvolvidos ao longo da implementa\c{c}\~ao.

\begin{lstlisting}[language=bash,caption={Criacao da diretoria de trabalho}]
mkdir -p /root/projeto
cd /root/projeto
\end{lstlisting}

A cria\c{c}\~ao desta diretoria permitiu centralizar os scripts do projeto num local \'unico, facilitando execu\c{c}\~ao, manuten\c{c}\~ao e testes.

\subsection{Implementa\c{c}\~ao do script}

O script \texttt{q1\_dns.sh} automatiza a cria\c{c}\~ao da zona \textit{master} DNS a partir do dom\'{\i}nio e do IP introduzidos pelo utilizador. O fluxo foi organizado em blocos: defini\c{c}\~oes iniciais, valida\c{c}\~oes, integra\c{c}\~ao no BIND, cria\c{c}\~ao da zona e valida\c{c}\~ao final.

\begin{lstlisting}[language=bash,caption={Definicoes base e validacoes iniciais}]
MAINCONF="/etc/named.conf"
CONFZONES="/etc/named/zones.conf"
ZONEDIR="/var/named"

read -p "Introduza o nome do dominio: " DOMINIO
read -p "Introduza o IP do servidor DNS/Web: " IP

if [ "$EUID" -ne 0 ]; then
    echo "Erro: execute como root."
    exit 1
fi
\end{lstlisting}

Este bloco define a base da solu\c{c}\~ao e valida a execu\c{c}\~ao como \texttt{root}, condi\c{c}\~ao necess\'aria para alterar ficheiros de sistema.

De seguida, s\~ao feitas valida\c{c}\~oes para garantir que dom\'{\i}nio, IP, pacotes e caminhos necess\'arios est\~ao dispon\'{\i}veis.

\begin{lstlisting}[language=bash,caption={Validacao de pre-condicoes de execucao}]
if [ -z "$DOMINIO" ] || [ -z "$IP" ]; then
    echo "Erro: dominio e IP sao obrigatorios."
    exit 1
fi

if ! command -v named-checkconf >/dev/null 2>&1; then
    echo "Erro: BIND/DNS nao esta instalado."
    echo "Execute: yum install bind bind-utils -y"
    exit 1
fi

if [ ! -f "$MAINCONF" ]; then
    echo "Erro: ficheiro $MAINCONF nao existe."
    exit 1
fi

if [ ! -d "$ZONEDIR" ]; then
    echo "Erro: diretoria $ZONEDIR nao existe."
    exit 1
fi
\end{lstlisting}

Estas verifica\c{c}\~oes evitam configura\c{c}\~oes incompletas e paragens por erro em ambiente n\~ao preparado.

Numa segunda fase, o script integra a nova zona no BIND atrav\'es do ficheiro \texttt{/etc/named/zones.conf}, evitando altera\c{c}\~oes repetidas no \texttt{named.conf}.

\begin{lstlisting}[language=bash,caption={Integracao da zona e prevencao de duplicados}]
touch "$CONFZONES"

if ! grep -q 'include "/etc/named/zones.conf";' "$MAINCONF"; then
    echo '' >> "$MAINCONF"
    echo 'include "/etc/named/zones.conf";' >> "$MAINCONF"
fi

if grep -q "zone \"$DOMINIO\"" "$CONFZONES"; then
    echo "Erro: a zona $DOMINIO ja existe em $CONFZONES"
    exit 1
fi

if [ -f "$ZONEFILE" ]; then
    echo "Erro: o ficheiro $ZONEFILE ja existe."
    exit 1
fi
\end{lstlisting}

Este bloco impede duplica\c{c}\~oes e torna o processo de cria\c{c}\~ao de zonas mais robusto.

Numa terceira fase, o script cria o ficheiro de zona com os registos m\'{\i}nimos (\texttt{SOA}, \texttt{NS} e \texttt{A}).

\begin{lstlisting}[language=bash,caption={Criacao do ficheiro de zona}]
cat > "$ZONEFILE" <<EOF
\$TTL 86400
@   IN  SOA ns1.$DOMINIO. root.$DOMINIO. (
        $SERIAL ; serial
        3600    ; refresh
        1800    ; retry
        604800  ; expire
        86400 ) ; minimum

@       IN  NS  ns1.$DOMINIO.
@       IN  A   $IP
ns1     IN  A   $IP
www     IN  A   $IP
EOF
\end{lstlisting}

Este \'e o n\'ucleo funcional do script, pois garante a gera\c{c}\~ao din\^amica da zona e o cumprimento do requisito de resolu\c{c}\~ao do dom\'{\i}nio para o IP indicado.

Ap\'os a cria\c{c}\~ao do ficheiro de zona, o script aplica permiss\~oes adequadas e regista formalmente a nova zona no ficheiro \texttt{zones.conf}.

\begin{lstlisting}[language=bash,caption={Permissoes e registo formal da nova zona}]
chown root:named "$ZONEFILE"
chmod 640 "$ZONEFILE"

cat >> "$CONFZONES" <<EOF

zone "$DOMINIO" IN {
    type master;
    file "$DOMINIO.zone";
    allow-update { none; };
};
EOF
\end{lstlisting}

A atribui\c{c}\~ao de permiss\~oes garante leitura correta pelo \texttt{named} e a entrada em \texttt{zones.conf} torna a zona ativa no BIND.

Por fim, o script valida toda a configura\c{c}\~ao antes de reiniciar o servi\c{c}o DNS.

\begin{lstlisting}[language=bash,caption={Validacao e ativacao da configuracao}]
named-checkconf || { echo "Erro no ficheiro $MAINCONF"; exit 1; }
named-checkzone "$DOMINIO" "$ZONEFILE" || { echo "Erro no ficheiro de zona"; exit 1; }
systemctl restart named
systemctl enable named
\end{lstlisting}

Este bloco confirma a validade da configura\c{c}\~ao antes de reiniciar o servi\c{c}o, reduzindo o risco de deixar o DNS em estado inconsistente.

\subsection{Execu\c{c}\~ao e demonstra\c{c}\~ao}

Ap\'os atribuir permiss\~ao de execu\c{c}\~ao ao script (\texttt{chmod +x /root/projeto/q1\_dns.sh}), a execu\c{c}\~ao foi conclu\'{\i}da com sucesso para o dom\'{\i}nio \texttt{teste.com}. A valida\c{c}\~ao confirmou a cria\c{c}\~ao do ficheiro de zona, o registo da zona em \texttt{/etc/named/zones.conf} e a resolu\c{c}\~ao correta no cliente.

\subsection{Evid\^encias visuais da Quest\~ao 1}

\begin{figure}[H]
\centering
\includegraphics[width=0.95\textwidth]{q1_execucao}
\caption{Execucao do script da Questao 1 e validacao da zona criada}
\label{fig:q1-execucao}
\end{figure}

\begin{figure}[H]
\centering
\includegraphics[width=0.95\textwidth]{q1_zonefile}
\caption{Conteudo do ficheiro de zona gerado automaticamente}
\label{fig:q1-zonefile}
\end{figure}

\begin{figure}[H]
\centering
\includegraphics[width=0.95\textwidth]{q1_zonesconf}
\caption{Definicao da zona adicionada em \texttt{/etc/named/zones.conf}}
\label{fig:q1-zonesconf}
\end{figure}

\begin{figure}[H]
\centering
\includegraphics[width=0.85\textwidth]{q1_nslookup_windows}
\caption{Validacao da resolucao DNS no cliente Windows}
\label{fig:q1-nslookup}
\end{figure}

\clearpage
\section{Quest\~ao 3: integra\c{c}\~ao DNS e VirtualHost Web}
\label{sec:q3}

\subsection{Objetivo}

A Quest\~ao 3 constitui uma extens\~ao direta da Quest\~ao 1. Para al\'em da cria\c{c}\~ao autom\'atica da zona DNS, o script passa tamb\'em a configurar automaticamente o servi\c{c}o Web, de forma a permitir que o dom\'{\i}nio criado responda por HTTP atrav\'es de um \textit{VirtualHost} dedicado no Apache.

Em termos pr\'aticos, a solu\c{c}\~ao implementada passa a abranger todo o ciclo de publica\c{c}\~ao de um dom\'{\i}nio, incluindo:

\begin{enumerate}
    \item cria\c{c}\~ao da zona DNS e do respetivo ficheiro de zona;
    \item cria\c{c}\~ao da diretoria Web associada ao dom\'{\i}nio;
    \item gera\c{c}\~ao autom\'atica do ficheiro de \textit{VirtualHost} no Apache;
    \item cria\c{c}\~ao de uma p\'agina inicial de boas-vindas;
    \item valida\c{c}\~ao e rein\'{\i}cio dos servi\c{c}os \texttt{named} e \texttt{httpd}.
\end{enumerate}

Com esta abordagem, o dom\'{\i}nio fica funcional de ponta a ponta: resolve em DNS e responde por HTTP com um \textit{VirtualHost} dedicado.

\subsection{Implementa\c{c}\~ao do script \texttt{q3\_dns\_web.sh}}

Numa primeira fase, o script define os caminhos principais utilizados tanto pelo BIND como pelo Apache.

\begin{lstlisting}[language=bash,caption={Variaveis principais para DNS e Apache}]
MAINCONF="/etc/named.conf"
CONFZONES="/etc/named/zones.conf"
ZONEDIR="/var/named"

HTTPD_CONF_DIR="/etc/httpd/conf.d"
WEBROOT_BASE="/var/www"
\end{lstlisting}

Estas vari\'aveis organizam a configura\c{c}\~ao de DNS e Web e tornam o script mais leg\'{\i}vel. Depois da recolha do dom\'{\i}nio e do IP, o script valida o ambiente, garante a inclus\~ao \'unica de \texttt{zones.conf} e impede duplica\c{c}\~oes de zona, ficheiro de zona e \textit{VirtualHost}.

Ap\'os estas valida\c{c}\~oes, o script cria a zona DNS e avan\c{c}a para a componente Web, gerando a diretoria do dom\'{\i}nio e a p\'agina inicial.

\begin{lstlisting}[language=bash,caption={Criacao automatica de WebRoot e pagina inicial}]
WEBROOT="$WEBROOT_BASE/${DOMINIO}"
mkdir -p "$WEBROOT"

cat > "$WEBROOT/index.html" <<EOF
<!DOCTYPE html>
<html lang="pt">
<head>
    <meta charset="UTF-8">
    <title>Bem-vindo a $DOMINIO</title>
</head>
<body>
    <h1>Bem-vindo ao dominio $DOMINIO</h1>
    <p>Este VirtualHost foi criado automaticamente.</p>
</body>
</html>
EOF
\end{lstlisting}

Este bloco prepara o conte\'udo Web m\'{\i}nimo do dom\'{\i}nio, permitindo validar imediatamente o funcionamento do \textit{VirtualHost}.

Depois de preparar a estrutura de conte\'udos, o script cria o ficheiro de \textit{VirtualHost} no Apache.

\begin{lstlisting}[language=bash,caption={Configuracao do VirtualHost no Apache}]
cat > "$VHOSTFILE" <<EOF
<VirtualHost *:80>
    ServerName $DOMINIO
    ServerAlias www.$DOMINIO
    DocumentRoot "$WEBROOT"

    <Directory "$WEBROOT">
        AllowOverride All
        Require all granted
    </Directory>

    ErrorLog /var/log/httpd/${DOMINIO}_error.log
    CustomLog /var/log/httpd/${DOMINIO}_access.log combined
</VirtualHost>
EOF
\end{lstlisting}

Este bloco integra o dom\'{\i}nio no Apache: \texttt{ServerName}/\texttt{ServerAlias} definem os nomes atendidos, \texttt{DocumentRoot} define o conte\'udo e os logs ficam separados por dom\'{\i}nio.

Por fim, o script valida as configura\c{c}\~oes dos dois servi\c{c}os antes de proceder \`a sua ativa\c{c}\~ao definitiva.

\begin{lstlisting}[language=bash,caption={Validacao final e reinicio dos servicos}]
named-checkconf || { echo "Erro no named.conf"; exit 1; }
named-checkzone "$DOMINIO" "$ZONEFILE" || { echo "Erro na zona DNS"; exit 1; }
httpd -t || { echo "Erro na configuracao Apache"; exit 1; }

systemctl restart named
systemctl restart httpd
\end{lstlisting}

Este bloco final valida DNS e Apache antes do rein\'{\i}cio dos servi\c{c}os, garantindo aplica\c{c}\~ao segura da configura\c{c}\~ao.

\subsection{Valida\c{c}\~ao funcional}

A valida\c{c}\~ao da Quest\~ao 3 foi feita em tr\^es n\'{\i}veis para confirmar o funcionamento das componentes DNS e Web.

\begin{enumerate}
    \item \textbf{Execu\c{c}\~ao do script no servidor}: confirma cria\c{c}\~ao do ficheiro de zona DNS, da entrada correspondente em \texttt{zones.conf}, do ficheiro de \textit{VirtualHost} no Apache, da diretoria Web do dom\'{\i}nio e da p\'agina inicial \texttt{index.html};
    \item \textbf{Resolu\c{c}\~ao DNS no cliente}: confirma que o dom\'{\i}nio configurado \'{e} corretamente resolvido para o endere\c{c}o IP do servidor;
    \item \textbf{Acesso HTTP no navegador}: confirma que o Apache entrega corretamente a p\'agina de boas-vindas criada automaticamente pelo script.
\end{enumerate}

\subsection{Evid\^encias visuais da Quest\~ao 3}

\begin{figure}[H]
\centering
\includegraphics[width=0.95\textwidth]{q3_execucao}
\caption{Execucao da Questao 3 no servidor}
\label{fig:q3-execucao}
\end{figure}

\begin{figure}[H]
\centering
\includegraphics[width=0.85\textwidth]{q3_nslookup_windows}
\caption{Validacao de resolucao DNS no cliente Windows}
\label{fig:q3-nslookup}
\end{figure}

\begin{figure}[H]
\centering
\includegraphics[width=0.95\textwidth]{q3_http_browser}
\caption{Acesso HTTP ao VirtualHost criado automaticamente}
\label{fig:q3-http}
\end{figure}
